\chapter{相关研究基础}
\label{chap:related-work}

本章从建筑幕墙巡检、视觉缺陷检测、AI Agent 工程范式和报告生成四个方向梳理本文工作的研究基础。通过相关研究分析可以看出，现有方法在单点检测能力上已经具有一定积累，但在项目级任务编排、Agent 受控调度和多源结果分层整合方面仍存在明显空白，这也构成本文技术方案设计的出发点。

\section{建筑幕墙巡检与数字化运维}
\label{sec:related-facade-inspection}

建筑幕墙巡检长期以来依赖人工现场勘查。人工巡检能够结合经验进行综合判断，但在高层建筑、大面积幕墙和复杂环境下效率较低，且存在安全风险。随着无人机技术和图像采集设备的发展，基于无人机的建筑立面巡检逐渐成为重要方向。Freimuth 和 K{\"o}nig 对无人机建筑检查的规划和执行流程进行了系统研究，说明无人机能够降低高空检测成本并提升数据采集效率\cite{freimuth2018uav}。在此基础上，红外热成像、图像配准和 BIM 结合等方法进一步扩展了建筑立面检测的数据来源和空间定位能力\cite{li2024facade,zhang2023uavbim}。

但是，数字化采集并不等于智能化评估。无人机可以快速采集图像，但图像采集之后仍需要图像管理、缺陷检测、结果解释和报告生成。若缺少后续软件系统，采集到的图像仍然是离散资料，难以转化为运维决策。因此，建筑幕墙评估系统需要将采集后的图像组织为项目资产，并建立从图像到评估报告的完整处理流程。

从工程应用角度看，幕墙巡检的难点并不只在“看清图像”。采集设备能够提高图像获取效率，检测模型能够提高局部缺陷识别效率，但真正面向业主、运维单位和工程师的交付物通常是带有结论、依据和建议的评估报告。这意味着系统需要同时处理视觉证据、结构化指标、项目语义和文本表达。若系统只停留在缺陷框或掩码图层面，仍然需要人工把这些结果重新解释、筛选、汇总和撰写。因此，图像检测之后的组织、解释和报告生成，是幕墙数字化运维走向智能化运维的关键环节。

此外，幕墙工程对象具有明显的项目属性。同一种缺陷出现在不同区域，其工程意义可能不同；同一张图像中的不同材料，也可能对应不同检测任务。传统视觉检测研究往往把图像视作独立样本，而实际运维更关心项目内图像之间的关系、图像组与建筑区域的关系、缺陷分布与维护优先级之间的关系。本文的系统设计正是围绕这种项目级应用需求展开。

\section{幕墙缺陷检测技术}
\label{sec:related-defect-detection}

深度学习推动了图像缺陷检测的发展。卷积神经网络能够自动学习图像特征，目标检测和分割模型能够输出缺陷区域、类别和置信度\cite{article1,conf1}。在建筑结构和幕墙场景中，裂缝检测、玻璃检测和表面缺陷识别等任务均已有较多研究。Ye 等人利用深度学习方法进行结构裂缝检测，为石材或混凝土裂缝识别提供了参考\cite{ye2019crack}；Mei 等人研究真实场景中的玻璃检测问题，对玻璃幕墙区域分析具有启发意义\cite{mei2020glass}。

从本文系统实现角度看，五类原子检测能力可以分别由 YOLO、传统计算机视觉算法或深度学习分类模型实现。金属锈蚀检测可以通过目标检测和分割获得锈蚀区域；石材裂缝检测可以通过分割方法获得裂缝掩码；石材污渍检测可以通过区域检测和局部分析得到污渍占比；玻璃平整度检测可以结合玻璃区域识别、边缘、直线、梯度和频域特征判断是否不平整；玻璃爆裂检测可以通过分类模型判断是否存在破损风险。

需要强调的是，与重新提出某一类检测模型的研究不同，本文更关注既有检测能力在工程系统中的组织、调用和结果整合。原因在于实际工程中，检测模型可以持续替换和迭代，系统真正需要解决的是“如何让这些检测能力被合理调用并形成评估报告”。因此，本文将检测模型抽象为原子能力，重点研究其工程封装、调度和结果整合。

这种定位并不削弱模型的重要性，而是重新定义模型在系统中的角色。对于工程软件而言，模型更像是一个可替换工具：当模型精度提高时，系统应能低成本接入；当新增缺陷类型时，系统应能扩展任务代码、数据表和展示逻辑；当某个模型执行失败时，系统应能隔离失败并保留其他任务结果。也就是说，模型能力需要被纳入工程治理，而不是以脚本形式孤立运行。本文后续提出的原子检测能力工具化，就是为了解决这一问题。

\section{从单点算法到项目级评估的差距}
\label{sec:related-algorithm-to-system}

已有幕墙缺陷检测研究通常以单一任务为目标，例如裂缝识别、玻璃检测或表面污渍检测。这类研究能够回答算法层面的问题：给定输入图像，模型是否能够准确输出目标区域。然而，项目级评估需要回答更复杂的问题：哪些图像需要执行哪些检测，多个检测结果之间如何互相补充，哪些缺陷对项目整体风险影响更大，报告中应如何表达证据与结论之间的关系。

从单点算法到项目级评估之间至少存在四个缺口。第一是数据组织缺口，图像需要按照项目、图像组、图像和检测任务多层组织。第二是任务编排缺口，系统需要根据图像内容动态选择检测能力，而不是机械执行全部模型。第三是结果解释缺口，模型输出的像素和坐标指标需要转化为工程语言。第四是报告生成缺口，系统需要把多张图像的局部结论融合为项目级评估意见。本文的研究价值正是在这些缺口处构建桥梁。

如果只关注检测模型，系统会缺少业务闭环；如果只关注 Web 管理界面，系统又缺少智能决策能力。本文选择的路线是以工程系统组织检测模型，以 Agent 承担语义决策和信息整合，从而把二者连接为完整应用。

\section{AI Agent 的发展与工程范式}
\label{sec:related-agent-paradigm}

大语言模型的兴起使 Agent 成为当前人工智能应用的重要范式。早期大语言模型主要用于文本生成和问答，而 Agent 更强调模型与外部工具、环境和用户目标之间的交互。ReAct 提出将推理和行动交替结合，使模型能够生成思考过程并执行任务动作\cite{yao2023react}。Toolformer 进一步说明语言模型可以学习何时调用工具、如何传递参数以及如何利用工具返回结果\cite{schick2023toolformer}。这些工作共同说明，模型本身的语言能力需要通过工具调用机制转化为可执行能力。

工业界和开源社区也逐渐形成 Agent 工作流开发范式。AutoGen 通过多智能体对话组织复杂任务，说明 Agent 可以被看作具有不同角色和能力的可编程组件\cite{wu2023autogen}。相关综述指出，大语言模型 Agent 可以应用于单智能体、多智能体和人机协作等场景，其关键能力包括规划、记忆、工具使用和反馈学习\cite{wang2023llmagents}。在工程实践中，Agent 往往不会完全自由行动，而是被放置在受控工作流中，通过结构化输入输出、工具白名单和状态管理保证系统可控。

本文正是基于这一认识，将 Agent 设计为业务链路中的受控智能节点。分类 Agent 不是自由选择任意工具，而是在预定义五类任务代码中选择；总结 Agent 不是直接修改系统状态，而是通过回调接口提交总结结果；项目报告 Agent 不是自行读取数据库，而是接收由 BIZ 服务组织好的项目上下文。这样的设计符合当前 Agent 工程化趋势：让模型承担语义理解和决策任务，让工程系统负责边界和状态。

\begin{table}[!htb]
  \caption{大语言模型应用范式的演进}\label{tab:llm-agent-evolution}
  \centering
  \begin{tabular}{@{}p{0.18\linewidth}p{0.35\linewidth}p{0.37\linewidth}@{}} \toprule
    \textbf{范式} & \textbf{主要特征} & \textbf{局限与本文启发} \\ \midrule
    文本生成 & 模型根据提示词直接输出自然语言。 & 适合问答和撰写，但难以与业务状态和外部工具形成闭环。 \\
    检索增强 & 先检索外部知识，再生成回答。 & 能补充上下文，但仍偏向知识问答，不直接解决任务调度问题。 \\
    工具调用 & 模型根据目标选择外部 API 或工具。 & 适合连接模型与工程能力，但需要严格约束工具集合和参数。 \\
    工作流 Agent & Agent 被放入确定业务流程，承担局部决策。 & 更适合工程落地，本文采用该路线实现检测路由和信息整合。 \\ \bottomrule
  \end{tabular}
\end{table}

\cref{tab:llm-agent-evolution} 展示了大语言模型应用范式从文本生成到工作流 Agent 的变化。本文将 Agent 理解为工程工作流中的智能决策节点，而非完全自主的软件机器人。这样的设计更符合垂直行业系统的要求：业务流程必须可解释，数据状态必须可追踪，外部工具调用必须可控，失败情况必须可恢复。

\section{Agent 工程化中的可控性问题}
\label{sec:related-agent-control}

Agent 系统虽然具备更强的自主性，但也带来工程可控性问题。大语言模型输出具有概率性，可能出现格式不稳定、内容越界、幻觉或工具选择错误等情况。对于普通问答系统，这些问题可能只是回答质量问题；对于工程业务系统，如果 Agent 输出直接驱动数据库状态或任务调度，则可能造成错误流程。因此，Agent 工程化必须关注结构化输出、工具白名单、状态隔离和人工复核。

本文采用的思路是将 Agent 决策限制在有限动作空间内。分类 Agent 只能输出五种任务代码的组合；图像总结 Agent 的结果只作为总结字段保存，不直接修改检测状态；项目报告 Agent 接收的输入由 BIZ 服务生成，输出也通过固定回调进入数据库。这样做牺牲了一部分开放性，但换来了工程系统所需的确定性和可追踪性。对于建筑幕墙评估这种专业场景，受控 Agent 比完全开放 Agent 更适合落地。

\section{工程软件中的异步任务编排}
\label{sec:related-async-workflow}

长耗时任务编排是智能检测系统必须解决的问题。模型推理、外部 Agent 调用、对象上传和报告生成都可能持续数秒到数分钟，若采用同步 HTTP 请求等待全部结果，系统容易超时，用户体验也较差。成熟工程系统通常采用任务队列、Worker、回调和状态查询机制，将“启动任务”和“获得结果”解耦。

本文系统采用类似思想，但结合 Agent 场景进行了扩展。用户启动检测后，BIZ 服务只负责创建任务并入队；分类、推理和总结能力由不同 Activity 服务异步完成；完成后通过回调上报结果；前端通过 WebSocket 感知状态变化。该模式使 Agent 和模型调用可以作为后台任务运行，也使系统能够在失败时记录状态并支持重试。

异步任务编排对本文尤为重要，因为 Agent 调用和模型推理都具有不可忽略的不确定性。Agent 平台可能因为网络、限流或输出异常导致失败；视觉模型可能因为图像格式、模型加载或运行环境导致失败；对象存储上传也可能受到网络影响。如果所有步骤都放在一次同步请求中，任何环节失败都会导致整个请求失败，且用户无法获得过程反馈。通过任务状态机和回调机制，系统能够把这些不确定因素转化为可记录、可展示和可恢复的状态变化。

从软件工程角度看，异步任务不仅是性能优化，也是复杂业务建模方式。检测主任务、子检测任务、总结任务和报告任务之间存在依赖关系。只有把这些关系显式建模为数据库状态和 Worker 编排逻辑，系统才能保证“分类失败不继续检测”“子任务全部成功才总结”“无检测任务时直接成功”等业务规则被稳定执行。

\section{报告生成与信息整合相关工作}
\label{sec:related-report-generation}

大语言模型擅长自然语言生成，但在知识密集型任务中，输入信息的组织方式直接影响输出质量。检索增强生成通过外部知识库向模型提供上下文，缓解模型仅依赖参数知识导致的不可靠问题\cite{lewis2020rag}。与此类似，工程报告生成也需要将外部数据组织成适合模型处理的上下文。如果输入过少，报告缺少依据；如果输入过多，模型容易遗漏重点或产生不稳定输出。

建筑幕墙韧性评估报告属于典型的多源信息整合任务。其输入包括项目基础信息、图像组信息、原子检测指标、图像产物、区域级数据、用户复核评论和任务状态。直接把这些信息一次性提交给大语言模型并不合理，因为底层指标数量可能随图像数量和区域数量快速增长。本文提出分层信息整合策略，将报告生成拆分为图像级和项目级两个阶段。该策略并非简单压缩文本，而是根据工程评估逻辑先形成单图语义，再进行项目级融合。

\section{本章小结}
\label{sec:related-gap}

综合来看，现有研究在三个方面尚存在不足。第一，幕墙检测研究多关注单类缺陷识别，对图像资产管理、任务调度和报告生成关注不足。第二，Agent 研究多集中在通用工具调用、代码执行和问答任务，在建筑工程垂直场景中的系统化落地仍有探索空间。第三，大语言模型报告生成往往关注生成效果本身，而对输入信息如何分层组织、如何与数据库状态和工程流程结合讨论不足。

本文的切入点正是在上述空白之间建立连接：以幕墙检测为现实场景，以 Agent 作为任务调度和信息整合能力，以工程系统作为稳定运行载体，构建从图像资产到韧性报告的完整链路。这样既避免了单纯算法研究难以落地的问题，也避免了单纯工程系统缺少智能亮点的问题。

更具体地说，本文有别于在单个模型指标上与已有算法竞争的路线，而是研究如何在真实业务流程中组织这些算法。已有检测模型回答“这张图是否存在某类缺陷”，本文系统进一步回答“这张图应该检测什么”“这些检测结果如何变成图像总结”“所有图像如何形成项目报告”。这三个问题正是 Agent 能力与工程链路结合的空间。

因此，本文的研究切入点具有复合性。一方面，它需要具备工程系统的完整性，能够支持登录、项目管理、图像上传、状态推送、结果展示和部署运行；另一方面，它需要体现 Agent 方法的必要性，即在任务路由和报告生成中解决传统规则系统难以处理的语义判断问题。本文后续章节将围绕这一主线展开：先定义问题和需求，再提出 Agent 调度与分层整合方法，最后说明工程架构如何支撑该方法落地。
